\clearpage
\VerdierPhysicalPageAnchor{0157}%
\begingroup\small\setlength{\parskip}{0pt}\setlength{\emergencystretch}{1.5em}
\VerdierSemanticAnchor{ch3-3-2-4-proof}
Pour démontrer le lemme (3.2.4), remarquons que le morphisme
$h(X')\xrightarrow{\operatorname{Im}((\operatorname{id}_{X'};f'))}
\operatorname{Im}(f')$ est un épimorphisme de $\mathcal D^\wedge$. Les
foncteurs représentables étant des objets projectifs de $\mathcal D^\wedge$,
il existe un morphisme $g\colon X\to X'$ tel
que (3.2.4.1) soit commutatif.

Passons maintenant à la construction du morphisme $g'\colon Y\to Y'$.
Il résulte des axiomes TRI et TRII des
catégories triangulées qu'il existe deux triangles distingués
$(Z,X,Y,u,f,v)$ et $(Z',X',Y',u',f',v')$. Comme, pour tout objet $W$ de
$\mathcal D$, le foncteur
$\operatorname{Hom}_{\mathcal D}(W,\mathord\cdot)$ est un foncteur
cohomologique (1.2.1), les
suites
\VerdierSemanticAnchor{ch3-3-2-4-exact-sequences}
\[
\begin{tikzcd}[column sep=3.7em]
h(Z)\arrow[r,"h(u)"]&h(X)\arrow[r,"\operatorname{Im}((\operatorname{id}_X;f))"]&\operatorname{Im}(f)\arrow[r]&0\\
h(Z')\arrow[r,"h(u')"']&h(X')\arrow[r,"\operatorname{Im}((\operatorname{id}_{X'};f'))"']&\operatorname{Im}(f')\arrow[r]&0
\end{tikzcd}
\]
sont exactes. Comme $h(Z)$ est un objet projectif de $\mathcal D^\wedge$, il
existe un morphisme $g''\colon Z\to Z'$ tel
que
\VerdierSemanticAnchor{ch3-3-2-4-triii-square}
\[
\begin{tikzcd}[column sep=5em,row sep=2.3em]
h(Z)\arrow[r,"h(u)"]\arrow[d,"h(g'')"']&h(X)\arrow[d,"h(g)"]\\
h(Z')\arrow[r,"h(u')"']&h(X')
\end{tikzcd}
\]
soit commutatif. Il résulte alors de l'axiome TRIII qu'il existe un morphisme
$g'$ tel que $(g'',g,g')$ soit un morphisme de triangle. En particulier, le
diagramme (3.2.4.2) est commutatif. Ceci achève la démonstration du lemme
(3.2.4) et de la proposition (3.2.3).

\VerdierSemanticAnchor{ch3-3-2-5}
\medskip\noindent\textbf{Proposition 3.2.5.}\quad
\textit{Soient $V$ et $V'$ deux objets admissibles de $\mathcal D^\wedge$ et
$\alpha\colon V\to V'$ un morphisme de $\mathcal D^\wedge$. Le noyau de
$\alpha$ dans $\mathcal D^\wedge$ (resp. le conoyau de $\alpha$ dans
$\mathcal D^\wedge$) est un objet admissible de $\mathcal D^\wedge$.}

\VerdierSemanticAnchor{ch3-3-2-5-setup}
D'après (3.2.3), on peut supposer que $V=\operatorname{Im}(f)$ et
$V'=\operatorname{Im}(f')$ et que $\alpha$ est le morphisme induit sur les
images des morphismes $h(f)$ et $h(f')$ par un diagramme commutatif :
\[
\begin{tikzcd}[column sep=5em,row sep=2.4em]
h(X)\arrow[r,"h(f)"]\arrow[d,"h(g)"']&h(Y)\arrow[d,"h(g')"]\\
h(X')\arrow[r,"h(f')"']&h(Y').
\end{tikzcd}
\]
\endgroup
